%% ---------------------------------------------------------------------------------------------------------
\section{Conclusion}\label{sec:conclusion}

A recovery intervention on a frozen policy can rescue failures and can also destroy successes, and the two should be measured separately. With paired branching from a common snapshot, an exact binomial certificate on a calibration split and a single held-out confirmatory test, an observation-only release-time centring controller was certified (exact tests, family-wise error 0.05, calibration-host configuration) to have a conditional harm rate below 5\% for two diffusion policies, and it raised their success from 87.5\% to 94.75\% and from 86.75\% to 93.25\% on new scenes. Much of the harm margin comes from abstention, and the abstention rules were not shown to be the cause of the low harm. The descriptive mechanism analysis follows a representative rescue and the single development harm through the controller as implemented, including the re-planning hand-back and the shortfall between commanded and executed motion, and the audit of the physical record states which measurements can be trusted. The protocol can be applied to other recovery interventions; whether this controller transfers to other tasks or to a real robot remains open.
